\documentclass[11pt]{article}

\usepackage[margin=1in]{geometry}
\usepackage{amsmath, amssymb}
\usepackage{parskip}
\usepackage{xcolor}
\usepackage[hidelinks]{hyperref}
\pagestyle{plain}

\newenvironment{reviewerquote}
  {\begin{quote}\itshape}
  {\end{quote}}
\newcommand{\pending}[1]{\textcolor{red}{\textbf{[To do: #1]}}}

\title{Response to reviewers\\
  \large Introspection Dynamics with Mutation in Additive Games}
\author{Harry Foster, Vince Knight, Sebastian Krapohl}
\date{\today}

\begin{document}

\maketitle

Dear Professor Zaccour,

We thank both reviewers for their careful reading of the manuscript and for
their constructive comments. We respond to each comment below; reviewer text
is set in italics, and our responses follow in upright type. Where a comment
has led to a change in the manuscript we describe the change.

\section*{Reviewer 1}

\subsection*{Main comments}

\begin{reviewerquote}
In the abstract, perhaps the introspection dynamics process could be better
explained to be more accessible to readers who are not familiar with it.
\end{reviewerquote}

We have rewritten the opening of the abstract to describe the process
explicitly: at each time step a randomly selected player compares their
current payoff to the payoff an alternative action would have given, and
switches with a probability increasing in the difference. Following a minor
comment below, we have also removed all formulas from the abstract.

\begin{reviewerquote}
I think the manuscript would sound stronger with a slight framing reordering:
the authors could first say they extend introspection dynamics to include
mutations, and from there they rederive results for additive games.
\end{reviewerquote}

Done. The abstract and the introduction now lead with the extension: we
first state that we extend introspection dynamics to include per-player,
per-action mutation and player-specific selection intensities, and then
present the additive-game results of Couto and Pal as rederived, and
generalised, from the extended dynamics.

\begin{reviewerquote}
The Introduction could include some more references. Think of a reader who
hasn't carefully read References [1] and [2]: it would be hard for them to
know how the present model relates to previous ones, such as, for example,
other EGT models for asymmetric games or logit-response dynamics.
\end{reviewerquote}

Done. The introduction now relates introspection dynamics to the
logit-response dynamics of evolutionary game theory in economics (Blume,
1993; Al\'os-Ferrer and Netzer, 2010), notes evolutionary approaches to
asymmetric games (McAvoy and Hauert, 2015), and adds the exploration
dynamics of Traulsen et al.\ (2009) for the role of mutation as
exploration, error, and behavioural noise.

\begin{reviewerquote}
Theorem 1's proof is intuitive (and perhaps even an obvious result). However,
it remains somewhat verbal. Perhaps the authors could complement it with a
more formal argument. For additive games, because players' updates are
independent of what others do, we can define a smaller state space for each
player, \{C, D\}. [...] It follows trivially that the stationary distribution
of such a transition matrix is \{p\_i, 1-p\_i\}.
\end{reviewerquote}

Done, following exactly the suggested argument. Section 3 now defines, for
each player, the selection kernel \(K_i\) on their own action set: the
stochastic matrix governing a selected player's next action, well defined
precisely when the game is additive. Theorem 1 verifies directly that the
product of the per-player stationary distributions is stationary for the
full chain, and that positive mutation makes the chain irreducible and
aperiodic, so the stationary distribution is unique. For two actions the two
rows of \(K_i\) coincide at \((p_i,\,1-p_i)\), which is the reviewer's
argument; the closed form for the marginal (Theorem 2) and the explicit
product measure (Theorem 3) follow. This per-player-chain framework also
supports the alternative mutation scheme raised by Reviewer 2.

\begin{reviewerquote}
Even though one cannot obtain closed formulas for non-additive games, it
would still be nice to see the effect of the mutations for those games too.
Would you observe similar patterns for non-linear PGGs? (E.g., can mutation
increase cooperation frequency when cooperation is not expected?)
\end{reviewerquote}

Done, in a new section (`Non-additive games') with a new six-panel figure.
We take the non-linear public goods game with synergy or discounting
parameter \(w\) (Hauert et al., 2006) as a case study and compute exact
quantities from the full \(2^N\) linear system, for parameters at which
defection dominates the linear game. The section reports the following.
First, the \(w=1\) values coincide with our closed form, a numerical check
of the additive theory. Second, with discounting the additive pattern
persists qualitatively: mutation raises low cooperation towards
\(\tfrac{1}{2}\), though no longer affinely. Third, and answering the
reviewer's question directly, with synergy mutation produces high
cooperation where none is expected: the game is bistable, the mutation-free
strong-selection dynamics started from defection never cooperate, yet any
positive mutation rate makes the chain ergodic and strong selection then
drives cooperation up towards \(1-\mu\) rather than down to \(\mu\).
Mutation acts as an equilibrium-selection device in non-additive games, and
we note the parallel with the role of rare mutations in stochastic
evolutionary dynamics. Fourth, beyond additivity there can be an optimal,
strictly positive mutation rate: for a bistable synergy game the exact
\(p_C\) is non-monotone in \(\mu\), peaking at \(\mu^*\approx0.02\), and
the optimum moves and its peak falls as selection strengthens; our
mutation-selection balance result rules this out for any additive game,
where \(p_C\) is affine in \(\mu\). Finally, we show that additivity is
exactly the point at which the players' actions decouple: heatmaps of the
pairwise correlations, for the same heterogeneous group as our validation
figure (contributions \(\alpha_i=i\)), show that every pairwise correlation
is identically zero at \(w=1\) despite the heterogeneity, weakly negative
under discounting, and positive and increasing in both players'
contributions when the two strict equilibria compete. Beyond additivity a
player's long-run behaviour is tied to their co-players' parameters, in
contrast to the additive case.

\begin{reviewerquote}
Reference [2] provides a closed formula for the stationary distribution for
additive games with any number of strategies. Would your extension with
mutations work similarly in the general case? Perhaps it would be worth
including such a result if possible (even in the simplest case of symmetric
mutations).
\end{reviewerquote}

Done, and for arbitrary player- and action-specific mutation probabilities
rather than only the symmetric case. We have gone further and restructured
the paper so that the general result is now the lead result: the setup of
Section 2 is stated for finite action sets \(A_i\), and Section 3 first
establishes that for any additive game in which player \(i\), when selected,
mutates to action \(b\) with probability \(\mu_{ib}>0\), the stationary
distribution of the full chain is the product of the per-player stationary
distributions (Theorem 1); the two-action closed form follows as a
specialisation (Theorems 2 and 3). Computing the stationary distribution
therefore reduces from one linear system of dimension \(\prod_i L_i\) to
\(N\) systems of dimension \(L_i\). We also record what is special to two
actions: only there do the rows of the per-player transition matrix
coincide, so for three or more actions the marginals are given by \(N\)
small linear systems rather than by a closed form.

\subsection*{Minor comments}

\begin{reviewerquote}
The manuscript is generally well-written, but it would benefit from another
proofreading, as there are a few typos and missing words/commas.
\end{reviewerquote}

We have made a proofreading pass alongside the corrections below, fixing in
particular a comma splice in the discussion of \(M_1\) and \(M_2\), one in
the proof of Theorem 1, and inconsistent equation referencing (all display
equations are now referenced with \verb|\eqref|). 

\begin{reviewerquote}
Introspection dynamics is generally defined for games with any number of
strategies. Therefore [...] it should read ``Introspection dynamics, in which
each player compares their payoff to their payoff under an alternative
action,'' [...] And specify that the linear system that needs to be solved to
obtain the stationary distribution has dimension \(2^N\) for the particular
case with 2 strategies.
\end{reviewerquote}

Done. The abstract and the introduction now say `an alternative action', and
the introduction states that the linear system has dimension \(2^N\), more
generally \(L^N\) when each player has \(L\) actions.

\begin{reviewerquote}
Abstract: I would personally avoid using formulas in the abstract, or else
all parameters should be defined (e.g., it's not clear what \(\mu\) is in the
abstract).
\end{reviewerquote}

We have removed all formulas from the abstract. The mutation probabilities
are now described in words (payoff-independent, player-specific probabilities
of switching to cooperation or to defection), and the closed form for the
public goods game is described as having \(N\) terms rather than \(2^N\).

\begin{reviewerquote}
Page 2: I'd write `general Prisoner's Dilemma' instead of `classical (RPST)
Prisoner's Dilemma'.
\end{reviewerquote}

Done, in both the introduction and Section 2.

\begin{reviewerquote}
Indentations, e.g., after Eq 1.
\end{reviewerquote}

Done. Text that continues a sentence after a display equation (after
equations (1) and (3)) is no longer indented.

\begin{reviewerquote}
Page 3: It may not be clear what \((1,1)\) means in
\(\Delta f_1((1,1))\) because that notation was not defined before.
\end{reviewerquote}

Done. The sentence now reads ``writing \((a_1, a_2)\) for the state in which
player 1 plays \(a_1\) and player 2 plays \(a_2\)'' before the notation is
used.

\begin{reviewerquote}
Page 3: ``Additivity for player 1 requires the equality \(T-R=P-S\)''.
Additivity actually requires this equality for all players. And perhaps it's
best to show the matrix form for the general Prisoner's Dilemma too. The
acronym `PD' has not been defined.
\end{reviewerquote}

Done. The text now reads ``Additivity requires, for each player, the
equality'', the general Prisoner's Dilemma is displayed in matrix form as
\(G_3\), and the acronym PD is defined at first use. (The example matrices,
previously \(M_1\), \(M_2\), \(M_3\), are now \(G_1\), \(G_2\), \(G_3\), to
avoid a collision with the total mutation probability \(M_i\).)

\begin{reviewerquote}
Page 3: ``(both differentials equal c).'' To be precise, it's \(c_1\) and
\(c_2\), unless you want to refer to the symmetric case.
\end{reviewerquote}

Done. The text now reads ``for player \(i\), both differentials equal
\(c_i\)''.

\begin{reviewerquote}
Page 3: ``To define introspection dynamics we introduce the following two
quantities:'' I wouldn't call the Fermi function and the mutation scheme
`quantities'. Also, using \(\mu_0\) for the probability of mutating to
cooperation and \(\mu_1\) for the probability of mutating to defection is a
counterintuitive notation given you denote cooperation as \(a=1\) and
defection as \(a=0\) before.
\end{reviewerquote}

We now write `components' rather than `quantities'. The mutation
probabilities have also been renamed throughout: they are now indexed by
the target action, \(\mu_{iC}\) for mutation towards cooperation and
\(\mu_{iD}\) for mutation towards defection, consistent with the general
notation \(\mu_{ib}\) for \(b\in A_i\); the notation \(\mu_{i0}\),
\(\mu_{i1}\) no longer appears.

\begin{reviewerquote}
Page 4: At times, you refer to C as a state. [...] To be precise, avoid
calling C a state. (An alternative would be defining another possible state
space for each player that is simply \{0,1\} or \{C, D\}.)
\end{reviewerquote}

Done. The wording now refers to selections `resulting in action \(C\)' rather
than placing a player `in state \(C\)', and, following the reviewer's
suggested alternative, each player's action set now serves as a formal
per-player state space: Section 3 defines the selection kernel \(K_i\) on
\(A_i\) (which is \(\{C,D\}\) in the two-action case) as part of the
restructured proofs (see our response to the main comments above).

\begin{reviewerquote}
Page 5: ``without mutation the chain degenerates to a point mass at the
dominant-action profile''. Consider complementing this observation with an
example to make it clearer for all readers.
\end{reviewerquote}

Done. We now give the donation game \(M_1\) as an example: each
\(\delta_i=c_i>0\), so as \(\beta\to\infty\) the mutation-free stationary
distribution collapses to a point mass on the all-defect profile.

\begin{reviewerquote}
Page 5: ``Mutation therefore extends the exact analysis to arbitrarily strong
selection.'' This sentence might also not be fully clear for readers not
familiar with introspection or Markov chains [...]
\end{reviewerquote}

Done. We have expanded this paragraph along the lines the reviewer sketches:
without mutation the chain loses ergodicity in the strong-selection limit, so
the limit must be taken after computing the stationary distribution at finite
selection intensity, as in [1]; mutation keeps the switching probabilities
bounded away from \(0\) and \(1\), so the chain stays ergodic at any
selection intensity and a unique stationary distribution always exists.

\begin{reviewerquote}
Page 8: In the second paragraph, you refer to the ``stationary cooperation
level of the linear public goods game in Proposition 3 of [2].'' Note that
the equation given in Proposition 3 of [2] is not exactly for the stationary
cooperation level, but rather the general stationary level for any state.
Also, this paragraph is just one sentence -- consider joining it with the
previous one.
\end{reviewerquote}

Done. The sentence now states that the product measure of Theorem 2 reduces
to the stationary distribution given, for every state, in Proposition 3 of
[2], with the cooperation level following; it has been joined to the
preceding paragraph.

\begin{reviewerquote}
Figure 1: In the legend, you mean Corollary 2 instead of 1. In the caption,
it should read ``All values were obtained using [3].''
\end{reviewerquote}

Done. The figure has been regenerated with the corrected legend, and the
caption now reads ``All values were obtained using [3].''

\begin{reviewerquote}
Table 1: ``Both columns to four decimal places''.
\end{reviewerquote}

Done. The caption now reads ``Both columns are given to four decimal
places.''

\begin{reviewerquote}
Page 9: It can be confusing to use the words `benefactor' and `beneficiary'.
Make it clearer what you mean by each.
\end{reviewerquote}

Done. We now define the terms once, immediately after Lemma 1: a player is a
\emph{net beneficiary} when \(r_i>N\), since each unit they contribute
returns more than a unit to them from the pool, and a \emph{net contributor}
when \(r_i<N\). These two terms are used consistently throughout; `benefactor'
and `non-beneficiary' no longer appear.

\begin{reviewerquote}
Remark 3: I believe it's not exactly a ``negative quantity \(\phi_i\)'', but
a negative sign stemming from \(\phi_i'\), to be more precise.
\end{reviewerquote}

Done. The remark now reads ``a negative sign stemming from the derivative
\(\phi_i'<0\)''.

\begin{reviewerquote}
Figure 2: The order in which each panel is mentioned in the text is switched.
The right panel is perhaps unnecessary because it's not showing any variation
(and that can be easily seen through the equations). Why not a plot that
varies over the mutation rates instead, for example?
\end{reviewerquote}

Done, on both counts. The right panel has been replaced by a plot of \(p_i\)
as a function of the symmetric mutation rate \(\mu\in[0,\tfrac{1}{2}]\) for
\(r_i<N\), \(r_i=N\), and \(r_i>N\), illustrating the affine
mutation-selection balance of Proposition 2, with all lines meeting at
\(p_i=\tfrac{1}{2}\) when \(\mu=\tfrac{1}{2}\). The surrounding text has been
reordered so that the panels are discussed left, centre, right.

\begin{reviewerquote}
Page 10: I personally prefer when a Conclusion section doesn't require the
reader to know all the notation introduced in the body text.
\end{reviewerquote}

Done. The conclusion has been rewritten so that, together with the
introduction, it is self-contained: the results are stated in words (the
players behave independently in the long run, each player's choice is
forgetful, and the computation reduces to one small linear system per
player) rather than through the notation of the body text.

\section*{Reviewer 2}

\begin{reviewerquote}
I'm worried that their result (Theorem 1) holds for a particular kind of
mutation model where players can mutate to both the actions (A or B)
independently, no matter which action they are playing currently. [...] It
turns out that if you only allow mutation to the other strategy, the
statement \(p(A\to B)=p(A\to A)=p_A\) (Theorem 1) is no longer true even if
the two mutational probabilities are the same. This means that Theorem 2 (the
factorization result) is not immediately true for this other (equally
reasonable or arguably more reasonable) mutational scheme. That said, Theorem
2 could hold regardless; I did not check that. Could the authors justify the
necessity of self-mutation (biologically or game theoretically)? I suggest
that the authors perform some analysis of this alternate mutational scheme
also.
\end{reviewerquote}

The reviewer is correct that under mutation restricted to the other action
the forgetful property of Theorem 1 fails: the probability that a selected
player's next action is \(C\) then differs, by the mutation rate, between the
two current actions. Our preliminary analysis shows that the factorisation of
Theorem 2 nevertheless survives: for an additive game each player still
evolves as an independent two-state Markov chain, whose stationary
distribution is explicit; with symmetric mutation rate \(\mu_i\) the marginal
becomes
\[
  p_i \;=\; \frac{(1-\mu_i)\,\phi_i(\delta_i)+\mu_i}{1+\mu_i}.
\]
\pending{Add a full section analysing this alternative mutation scheme
(including action-dependent rates), proving the modified marginals and the
persistence of the product measure, and comparing the two schemes. Add a
justification of the original scheme: drawing the mutant action from the full
action set, including the current one, is the standard convention for
mutation and exploration in evolutionary and learning dynamics, and we will
cite the relevant literature.}

\vspace{2em}

Thank you again for the constructive reports.

Harry Foster, Vince Knight, and Sebastian Krapohl

\end{document}
